% =====================================================================
% Section 5 — Discussion
%
% Same provenance rule as Section 4: every quantitative statement here
% resolves to a table or figure in Section 4, which resolves to registry
% 20260819T043733Z. This section adds interpretation and limitations; it
% introduces no new numbers.
%
% Dispersion is always a standard deviation over three seeds, never an
% interval estimate.
% =====================================================================

\section{Discussion}
\label{sec:discussion}

\subsection{Architecture Ranking and Seed-to-Seed Variation}
\label{sec:discussion:ranking}

The benchmark separates into two findings of very different strength, and
they should not be reported with the same confidence.

The robust finding is that ResNet-50 is not competitive on this task. Its
$92.356 \pm 1.022$\,\% accuracy sits roughly seven points below
ConvNeXt-Tiny, a gap several times its own seed-to-seed standard deviation,
and it produces 122 of the 204 errors made by all five backbones across all
seeds (Section~\ref{sec:results:perspecies}). The gap is also structural
rather than diffuse: it is concentrated in specific species where the other
four architectures are near-perfect.

The fragile finding is the ordering among the remaining four. ConvNeXt-Tiny
leads Swin-Tiny by 0.438 accuracy points while their standard deviations are
0.354 and 0.581; MobileNetV3-Large trails Swin-Tiny by 0.752 points with a
standard deviation of 0.319. These separations are of the same order as the
seed-to-seed variation, so with three seeds the ordering cannot be treated
as established, and a different seed triple could plausibly reorder them.
The defensible claim is that ConvNeXt-Tiny, Swin-Tiny and MobileNetV3-Large
all reach a comparable operating point, and that identifying a single best
architecture among them would require more seeds than this study ran.

The MobileNetV3-Large result deserves emphasis for a different reason. It
reaches within 1.19 accuracy points of the best model using 4.22\,M
parameters against 27.83\,M, and 0.22\,G MACs against 4.46\,G
(Table~\ref{tab:efficiency}) --- roughly one sixth the parameters and one
twentieth the compute. For the field-deployment scenario that motivates
macroscopic wood identification, that trade-off matters more than the
ranking at the top. The comparison is one of accuracy against static
complexity only; inference latency is not reported here, for the reason
given in Section~\ref{sec:results:efficiency}.

\subsection{What the Models Are Using: Evidence Against an Anatomical Reading}
\label{sec:discussion:shortcut}

The most consequential limitation of this benchmark is not the size of the
accuracy gaps but what the accuracies measure. The probes in
Section~\ref{sec:results:convergence} show that a Swin-Tiny trained on
$16\times16$ inputs reaches 94.361\,\% accuracy, and one trained on heavily
blurred inputs 95.677\,\%, against 99.060\,\% for the full-resolution
reference. Both destroy the anatomical structure that macroscopic wood
identification relies on, and both retain more than 95\,\% of full-image
performance.

The natural first explanation --- that the models are simply reading
colour --- is not supported: the histogram-only probe reaches 37.970\,\%.
Colour carries real signal, far above the 5.882\,\% chance level, but it
accounts for well under half of what the spatial probes achieve. The
decisive evidence is therefore spatial but low-frequency: something that
survives sixteen-fold downsampling. Broad tonal gradients, coarse
figure, overall lightness distribution and specimen-level acquisition
characteristics are all consistent with this and are not separated by the
present experiments.

Two readings remain open, and they have very different consequences. The
benign reading is that conifer species genuinely differ in bulk optical
properties at this scale, in which case the models have learned a real but
coarser cue than intended, and the benchmark measures a legitimate if
narrower task. The adverse reading is that low-frequency content carries
specimen identity --- lighting, staining, mounting, camera setup --- and
that residual specimen leakage (Section~\ref{sec:discussion:originkey})
allows it to be exploited. The two readings cannot be separated from the
present data. Distinguishing them requires evaluation on independently
acquired images, which this study does not have
(Section~\ref{sec:discussion:external}).

The consequence for reporting is direct. The headline accuracies should be
read as accuracy on this task and this corpus, not as evidence that the
models identify wood by anatomy, and the project's own probe protocol
records the benchmark as not yet released on this basis. This does not
invalidate the architecture comparison in
Section~\ref{sec:discussion:ranking}: all five backbones were trained and
evaluated under identical conditions, so their relative ordering remains a
valid comparison. It restricts what the absolute numbers may be claimed to
demonstrate.

Two further observations follow from the probes. Their retention ordering
(blur above downsampling) suggests the useful content is at a coarser scale
than $16\times16$ sampling preserves, which is a testable prediction for a
resolution sweep --- an experiment that was configured for this study but
never run (Section~\ref{sec:discussion:scope}). And the probes were trained
for 20 epochs against 50 for the main runs, so their absolute numbers are
lower bounds on what the destroyed inputs can support, which strengthens
rather than weakens the argument.

\subsection{Rotation Does Not Improve Uniformly}
\label{sec:discussion:rotation}

Discrete $90^\circ$ rotation was motivated by the observation that a
macroscopic cross-section has no canonical orientation, so rotations by
multiples of $90^\circ$ should preserve every diagnostic feature. The data
do not support the expected benefit.

The effect changes sign across architectures: positive for ResNet-50 and
EfficientNet-B3, exactly zero for MobileNetV3-Large, negative for both
ConvNeXt-Tiny and Swin-Tiny (Table~\ref{tab:ablation}). An augmentation that
genuinely encoded a property of the imaging geometry would be expected to
help, or at least not hurt, regardless of backbone. That the two strongest
models are the two that decline is the opposite of the pattern a real
invariance gain would produce.

Two readings remain open and this study cannot separate them. The
augmentation may be redundant --- the main configuration already applies
continuous $\pm30^\circ$ rotation plus horizontal and vertical flips, so
much of the orientation invariance may already be covered, leaving discrete
$90^\circ$ rotation to trade away sample diversity for nothing. Or the
deltas may simply be seed noise: four of the five magnitudes fall within
about 1.6 times the corresponding seed-to-seed standard deviation, and with
a single seed per ablation there is no way to tell a small real effect from
a favourable draw. Either way, the practical conclusion is the same:
discrete $90^\circ$ rotation should not be adopted as a default for this
task on the strength of these runs.

Label smoothing behaves differently and is worth separating. Removing it
lowers accuracy for all five backbones without exception. No individual
delta is large, and three of five are below the corresponding seed-to-seed
standard deviation, so the effect size is not established --- but the
consistency of direction across five independent architectures is the one
pattern in the ablation table that does not look like noise. Raising label
smoothing to 0.2 shows no such consistency. This is weak evidence that the
default of 0.1 is a reasonable choice, and it is not evidence that 0.1 is
optimal.

\subsection{The Single-Seed Ablation Design Is the Binding Limitation}
\label{sec:discussion:oneseed}

Every ablation in this study was trained with seed 42 only. This is the most
consequential methodological limitation of the ablation results, and it is
not repairable by analysis.

The main runs establish that seed-to-seed standard deviation on this task
ranges from 0.319 to 1.022 accuracy points depending on backbone
(Table~\ref{tab:task_b}). Of the 15 ablation deltas in
Table~\ref{tab:ablation}, 9 are smaller than the standard deviation of their
own backbone, and only one exceeds it by more than about 2.4 times. A
single-seed delta of a fraction of a point therefore carries almost no
information about the underlying effect: the measurement noise and the
effect being measured are the same size.

This has a specific consequence for how the ablation table should be read.
Individual cells should not be interpreted. What can be read is the pattern
across the row of five backbones --- a consistent sign across five
independently initialised architectures, as seen for label smoothing 0.0, is
harder to explain by seed noise than any single cell, even though each cell
alone is uninformative. This is why Section~\ref{sec:results:ablation}
reports directional consistency rather than effect magnitudes, and why no
ablation is presented as an improvement.

Repeating each ablation across the same three seeds used for the main runs
would cost 30 additional training runs and would make these comparisons
interpretable. That is the single highest-value extension to this study.

\subsection{Effect of Origin-Key Construction}
\label{sec:discussion:originkey}

The Task~B partition groups images by specimen so that all crops from one
physical wood sample fall in a single split. This matters because the source
corpus is heavily cropped: 5{,}830 images derive from 557 specimens, and
88.3\,\% of the images come from specimens contributing 24 crops or more.
Under an image-level split, a specimen with 30 crops has a probability of
about $6.9\times10^{-16}$ of avoiding the training set entirely, so
essentially every large specimen appears in training, and test images become
different crops of familiar samples rather than new samples. The earlier
image-level Task~B partition showed exactly this: 97.62\,\% of its test
images came from specimens already present in training. Accuracy measured
under that partition describes crop recognition, not species generalisation.

The specimen grouping used here removes that particular failure, but the
guarantee it provides is bounded by how the specimen key is constructed. The
key is derived from the filename: the trailing crop index is stripped and the
class name prefixed, so \texttt{3351.\ Araucaria araucana.1.(12).png} becomes
\texttt{Araucaria araucana::3351.\ Araucaria araucana.1.}. It is a naming
convention, not a specimen identifier recorded at collection time. Two
distinct consequences follow, and they are not symmetric.

Over-grouping --- two physically different specimens sharing a filename stem
and being merged into one key --- costs only split flexibility. It is
conservative and does not inflate results.

Under-grouping is the real risk. If crops of one physical specimen were
saved under different stems, for example across separate scanning sessions or
after a renaming step, the key assigns them different origins and the
grouping silently fails to separate them. Nothing in the filename can detect
this, and no independent specimen identifier exists in the corpus against
which to check it.

A dedicated audit of the key gives this concern a magnitude. Regrouping the
552 keys of the 17-class set under a coarser candidate rule yields 325
groups, a 41.1\,\% collapse: on that reading, four keys in ten are fragments
of a larger physical unit rather than units in their own right. The audit
also finds 86 sample numbers reused across species, so the numeric part of a
filename is not by itself a specimen identifier. Its stated conclusion is
that whether the key denotes a physical specimen or merely an image or view
cannot yet be determined, and that filename structure must not substitute
for verification against the collection's own records; that verification
remains outstanding. The strength of the specimen-separation claim therefore
rests on the source collection's naming discipline, which this study did not
independently verify.

Two further caveats affect what can be checked from the published artifacts.
First, roughly 19\,\% of test images carry renumbered filenames in the
deployed split, so their origin key cannot be reconstructed after the fact;
among the 81\,\% that can be reconstructed, no specimen is shared between
validation and test. Second, the registry publishes predictions for the
validation and test splits only, so the train-versus-test separation --- the
one that actually matters --- cannot be re-verified from the released
artifacts and rests on the split construction itself. Publishing the
specimen-to-split assignment alongside the registry would close this gap.

\subsection{No External or Field Validation}
\label{sec:discussion:external}

All results come from a single corpus, split internally. The models have
never been evaluated on wood imaged by a different operator, with different
equipment, under different lighting, or at a different site. Accuracy above
99\,\% should be read as accuracy on held-out specimens \emph{from this
collection}, not as an estimate of field performance.

This limit is more restrictive than it may appear, for a reason specific to
this study. The single confusion that survives in the strongest model is
between two \textit{Pinus} species, and the errors that separate ResNet-50
from the rest are concentrated in a handful of species. Both patterns are
properties of this corpus at its current size --- 24 to 39 test images per
species, drawn from a limited number of specimens per species. A different
collection could easily show a different set of difficult species without
contradicting anything reported here.

Nothing in this study supports a claim about deployment readiness. The
appropriate next step is evaluation on an independently collected set of
specimens, ideally imaged under conditions the models have not seen, before
any operational use is considered.

\subsection{Scope Decisions: Task A and the Resolution Ablation}
\label{sec:discussion:scope}

Two experiments that appear in the project plan are not reported in
Section~\ref{sec:results}, and the reasons differ.

\paragraph{Task A (family-level classification).} The family-level task used
a species-stratified split operating at the image level, the same mechanism
that produced 97.62\,\% specimen coverage in the earlier Task~B partition
(Section~\ref{sec:discussion:originkey}), applied to the same heavily cropped
corpus. Its results are therefore subject to the same leakage, and its
near-perfect family-level scores --- including one architecture at
100.000\,\% on all three seeds --- cannot be read as evidence of
generalisation to unseen specimens. Reporting Task~A alongside a
specimen-grouped Task~B would place two incomparable evaluation protocols in
one results section and invite exactly the misreading the grouped split was
built to prevent. Task~A is therefore excluded rather than reported with a
caveat. Restoring it would require rebuilding the family-level split by
specimen and retraining all 15 runs; at family level this is less likely to
run short of classes than at species level, since the three species excluded
from Task~B belong to families that retain many other specimens. Note also
that the family-level task is defined over four families, folding
\textit{Araucaria araucana} into Pinaceae, whereas the grouping used in
Section~\ref{sec:results} keeps Araucariaceae separate; the two would have to
be reconciled before the tasks could be reported side by side.

\paragraph{Input resolution.} A $384$-pixel resolution ablation was
configured but never trained; no runs for it exist in the registry. Input
resolution is consequently untested in this study. This is a genuine gap
rather than an exclusion: resolution plausibly interacts with the one
difficulty that survives in the best model, since distinguishing
\textit{P.\ krempfii} from \textit{P.\ dalatensis} may depend on fine
anatomical detail that $224$-pixel inputs do not resolve. Testing it is a
natural companion to the multi-seed ablation proposed in
Section~\ref{sec:discussion:oneseed}.
